% === Begin Label Data ===
% ID: 829
% 难度星级: 
% 标签: 
% 备注: 
% 组卷引用次数: 0
% === End  Label Data ===

\begin{problem}{2012}{G}{重庆卷（理）}{15}{排列组合}
某艺校在一天的6节课中随机安排语文、数学、外语三门文化课和其他三门艺术课各1节，则在课表上的相邻两节文化课之间最多间隔1节艺术课的排法共有\underline{\hspace{4em}}种（用数字作答）.
\end{problem}

\begin{answer}
$\dfrac{3}{5}$
\end{answer}

\begin{solutions}
分三类：第一类：每相邻两节文化课之间都间隔1节艺术课，有 $ 2P_3^3P_3^3 $ 种排法；

第二类：三节文化课相邻，有 $ P_3^3P_4^4 $ 种排法；

第三类：三节文化课中有两节相邻，与另一节之间间隔1节艺术课，有 $ C_3^2P_2^2C_1^1P_2^2P_3^3 $ 种排法。

故共有 $ 2P_3^3P_3^3 + P_3^3P_4^4 + C_3^2P_2^2C_1^1P_2^2P_3^3 = 432 $ 种排法，所以所求概率为 $ P = \dfrac{432}{P_6^6} = \dfrac{3}{5} $.
\end{solutions}